\chapter{Thực nghiệm và đánh giá}

\section{Bộ dữ liệu}
Mô tả bộ dữ liệu được dùng trong thí nghiệm, bao gồm nguồn gốc, số lượng mẫu, kiểu nhãn, điều kiện thu thập và các tiêu chí tiền xử lý dữ liệu đầu vào. Nếu dữ liệu được chia theo nhiều tập con, cần nêu rõ cách chia và lý do chọn tỉ lệ đó.

\section{Thiết lập thực nghiệm}
\subsection{Phần cứng và phần mềm}
Ghi rõ cấu hình máy, môi trường chạy, phiên bản ngôn ngữ lập trình và các thư viện chính được sử dụng trong báo cáo.

\subsection{Siêu tham số và quy trình huấn luyện}
Trình bày batch size, learning rate, số epoch, kích thước ảnh, chiến lược tối ưu, tiêu chí dừng sớm và các tham số đáng chú ý khác.

\subsection{Kịch bản thí nghiệm}
Mỗi thí nghiệm nên chỉ ra rõ baseline, phương pháp cải tiến và điều kiện so sánh để đảm bảo tính công bằng.

\section{Kết quả định lượng}
\subsection{Bảng so sánh các phương pháp}
Kết quả cần được trình bày dưới dạng bảng để so sánh giữa các mô hình hoặc các cấu hình khác nhau.

\begin{table}[H]
    \centering
    \caption{Mẫu bảng so sánh kết quả thực nghiệm.}
    \label{tab:results_template}
    \begin{tabular}{|l|c|c|c|l|}
        \hline
        \textbf{Phương pháp} & \textbf{Chỉ số 1} & \textbf{Chỉ số 2} & \textbf{Runtime} & \textbf{Ghi chú} \\
        \hline
        Baseline & ... & ... & ... & ... \\
        \hline
        Cải tiến 1 & ... & ... & ... & ... \\
        \hline
        Cải tiến 2 & ... & ... & ... & ... \\
        \hline
    \end{tabular}
\end{table}

\subsection{So sánh baseline và phương pháp cải tiến}
Phần này cần nêu rõ phương pháp nào tốt hơn, tốt hơn ở chỉ số nào và đánh đổi gì về tốc độ hoặc độ phức tạp.

\subsection{Kết quả theo từng trường hợp}
Nên bổ sung các ví dụ thành công, ví dụ thất bại và các trường hợp biên để người đọc thấy rõ hành vi của mô hình.

\section{Phân tích lỗi}
\subsection{Trường hợp tốt}
Chỉ ra các mẫu mà hệ thống dự đoán chính xác và lý do vì sao mô hình hoạt động tốt trong các trường hợp này.

\subsection{Trường hợp xấu}
Mô tả các trường hợp hệ thống dự đoán sai, ví dụ do nền phức tạp, vật thể nhỏ, ánh sáng kém hoặc nhãn mơ hồ.

\subsection{Nguyên nhân và hạn chế}
Phân tích nguyên nhân gốc của lỗi thay vì chỉ liệt kê kết quả. Có thể liên quan đến dữ liệu, tiền xử lý, mô hình hoặc cách đánh giá.

\section{Thảo luận}
Tổng hợp lại ý nghĩa của các kết quả, mức độ phù hợp của phương pháp với bài toán thực tế và những điểm cần cải thiện nếu triển khai trong môi trường thật.